\section{Isolated stars of the initial mesh}
\label{sec:area_law_initial_stars}

We use the normalized fan of Definition~\ref{def:area_law_repair_fan},
with its eight triangles $T_s$, $s\in J$, the final endpoint $b_s$ of
segment $s$, and the counterclockwise successor $\tau$ of
Definition~\ref{def:area_law_cell_fan_successor}. For $v\in\mathbb R^2$
and $\rho>0$ put $W_s(v,\rho)=\{v+\rho x:x\in T_s\}$ and
$Q_\rho(v)=v+[-\rho,\rho]^2$. Distances are maximum-norm distances.
The endpoints $b_s$ are the eight nonzero vectors with coordinates in
$\{-1,0,1\}$, that is, the axis and diagonal directions.

% Source: 10-geometry.tex, geometry:initial-stars, lines 332--349;
% geometry:repair-induction, lines 471--485.
% Revision: adc7f1241b42e322a6451854ab7e4b4c146bf78a.
\begin{definition}[Isolated star]
    \label{def:area_law_isolated_star}
    \lean{TNLean.PEPS.AreaLaw.Geometry.repairFanDirection,
        TNLean.PEPS.AreaLaw.Geometry.repairFanDirection_mem,
        TNLean.PEPS.AreaLaw.Geometry.IsIsolatedStar,
        TNLean.PEPS.AreaLaw.Geometry.RepairStage.HasIsolatedStars,
        TNLean.PEPS.AreaLaw.Geometry.InitialIsolatedStarsStatement}
    \leanok
    \uses{def:area_law_repair_fan,def:area_law_repair_current_assignment,
        def:area_law_cell_fan_successor,def:area_law_repair_stage}
    Let $(R_i)_{i\in I}$ be open regions with colors
    $\xi:I\to\{0,1\}$. The square $Q_\rho(v)$ is an isolated star of
    this family when there is $\sigma:J\to I$ with the following
    properties.
    \begin{enumerate}
        \item For every $s$, $\operatorname{int}W_s(v,\rho)\subseteq
                  R_{\sigma(s)}$.
        \item For every $i$,
              $\overline{R_i}\cap Q_\rho(v)
                  =\bigcup_{\sigma(s)=i}W_s(v,\rho)$.
        \item For every $x\in Q_\rho(v)$, the point $x$ lies on some
              frontier $\partial R_i$ if and only if $x-v$ is a
              nonnegative multiple of $b_s$ for some $s$ with
              $\sigma(s)\ne\sigma(\tau(s))$.
        \item For every $s$, $\xi(\sigma(s))=\xi(\sigma(\tau(s)))$ if
              and only if $\sigma(s)=\sigma(\tau(s))$.
        \item If $\sigma(s)=\sigma(\tau(s))$ for every $s$, then
              $Q_\rho(v)\subseteq R_i$ for some $i$.
    \end{enumerate}
    The rays in the third property are the active rays. A run of
    triangles between consecutive active rays has one identifier, so it
    lies in one region together with its seams. A stage
    $\mathcal D$ has isolated stars with constant $a_0$ when
    $\operatorname{dist}(v_p,v_q)\ge a_0\max(S_p,S_q)$ for distinct
    occurrences $p\ne q$, and $Q_{a_0S_p}(v_p)$ is an isolated star of
    the current regions for every occurrence $p$. The initial-star
    assertion for the recursive family states that there is $a_0>0$ such
    that, for every width $C_0$, the initial stage has isolated stars
    with constant $a_0$ whenever $k_0$ is large enough.
\end{definition}

% Source: 10-geometry.tex, geometry:initial-stars, lines 340--349;
% geometry:repair-induction, lines 486--494.
\begin{lemma}[Shrinking an isolated star]
    \label{lem:area_law_isolated_star_mono}
    \lean{TNLean.PEPS.AreaLaw.Geometry.repairCurrentFanTriangle_inter_closedBall,
        TNLean.PEPS.AreaLaw.Geometry.IsIsolatedStar.mono,
        TNLean.PEPS.AreaLaw.Geometry.RepairStage.HasIsolatedStars.mono}
    \leanok
    \uses{def:area_law_isolated_star}
    For $0<\rho'\le\rho$ one has
    $W_s(v,\rho)\cap Q_{\rho'}(v)=W_s(v,\rho')$. Hence an isolated star
    $Q_\rho(v)$ gives an isolated star $Q_{\rho'}(v)$ with the same
    assignment. If every scale of a stage is positive, isolated stars
    with constant $a_0$ give isolated stars with every constant
    $a_0'\in(0,a_0]$.
\end{lemma}
\begin{proof}\leanok
    A point of $T_s$ is $ux$ with $u\in[0,1]$ and $x$ on the outer side,
    where $|x|=1$. Thus $v+\rho ux$ lies in $Q_{\rho'}(v)$ exactly when
    $\rho u\le\rho'$, that is, when it equals $v+\rho'u'x$ with
    $u'=\rho u/\rho'\in[0,1]$. The five properties restrict to the
    smaller square by this identity.
\end{proof}
